% ============================================================================
% Documentación: Búsqueda Tabú Reactiva (RTS) para CARP
% Generado/actualizado: 2026-07-10. Fuente: metacarp/busqueda_tabu_reactiva.py,
% scripts/run_tabu_reactiva_automatico.py, _calibracion_2knob.
% ============================================================================
\section{Búsqueda Tabú Reactiva (RTS)}
\label{sec:mh-rts}

\subsection{Concepto algorítmico}
La RTS (Battiti y Tecchiolli, 1994) extiende la Búsqueda Tabú clásica haciendo
que la lista tabú \textbf{deje de ser estática}: el algoritmo detecta cuándo
está ciclando y ajusta su memoria en consecuencia. Sus tres pilares:

\begin{enumerate}
  \item \textbf{Tenencia tabú dinámica}: el tenure varía dentro de
    $[\texttt{tenure}_{\min}, \texttt{tenure}_{\max}]$. Cuando la solución
    actual ya fue visitada antes (señal de ciclo), el tenure crece
    multiplicándose por \texttt{factor\_aumento} $(>1)$; tras
    suficientes iteraciones sin repeticiones decrece multiplicándose por
    \texttt{factor\_reduccion} $(<1)$.
  \item \textbf{Memoria de soluciones visitadas}: cada solución se registra
    mediante un hash canónico en un diccionario, lo que permite detectar
    repeticiones en $O(1)$ sin comparar contra todo el historial.
  \item \textbf{Mecanismo de escape}: si una solución se repite más de
    \texttt{umbral\_repeticiones\_escape} veces (ciclo ``duro''), se aplican
    varios movimientos aleatorios consecutivos ignorando la lista tabú, y se
    limpian tanto la lista como el historial (las prohibiciones de la zona
    vieja no son útiles en la nueva).
\end{enumerate}

Hereda de la TS simple el muestreo de lote (\texttt{tam\_vecindario} vecinos,
evaluación vectorizada), el criterio de aspiración, el doble criterio de
parada y los 9 operadores compartidos del proyecto.

\subsection{Parámetros}
\begin{table}[htbp]
\centering
\caption{Parámetros de la RTS.}
\begin{tabular}{lll}
\toprule
Parámetro & Significado & Valor calibrado \\
\midrule
\texttt{factor\_aumento} & multiplicador del tenure al detectar ciclo & 1.2 (grid) \\
\texttt{factor\_reduccion} & multiplicador de relajación del tenure & 0.95 (knob 2) \\
\texttt{umbral\_repeticiones\_escape} & repeticiones antes del escape & grid \\
$p_{inter}$ & prob.\ de operadores inter-ruta (factible) & 0.5 \\
$[\texttt{tenure}_{\min}, \texttt{tenure}_{\max}]$ & cotas del tenure dinámico & instance-aware \\
\bottomrule
\end{tabular}
\end{table}

\subsection{Búsqueda de malla (grid search)}
Es la grilla más grande de la campaña --- la RTS es la metaheurística con más
knobs estructurales:
\begin{itemize}
  \item \textbf{Grilla principal} (\texttt{run\_tabu\_reactiva\_automatico.py}):
    \texttt{factor\_aumento} $\in \{1.05, 1.1, 1.2, 1.3, 1.4\}$ (5) $\times$
    \texttt{umbral\_escape} $\in \{2, 3, 5, 8\}$ (4) $\times$
    $p_{inter} \in \{0.4, 0.5, 0.6, 0.7, 0.8\}$ (5) $\times$
    \texttt{factor\_reduccion} $\in \{0.85, 0.9, 0.95\}$ (3)
    $= 300$ combinaciones $\times$ 23 instancias $\times$ 5 repeticiones
    $= 34{,}500$ corridas.
  \item \textbf{Segundo knob} (verificación fina): \texttt{factor\_reduccion}
    $\in \{0.85, 0.9, 0.95\}$ con el resto fijo; ganó $0.95$ con gap medio de
    $5.101\,\%$ (contra $5.438\,\%$ y $5.766\,\%$). Reducción lenta del tenure
    = memoria que se relaja con cautela.
  \item \textbf{Configuración fija resultante}: \texttt{factor\_aumento} $=
    1.2$, \texttt{factor\_reduccion} $= 0.95$, $p_{inter} = 0.5$; corrida con
    los tres approaches a 5 repeticiones.
\end{itemize}

\paragraph{Resultado en las 23 pequeñas.} Mejor-por-instancia: gap medio
$1.16\,\%$ y 16/23 BKS --- segunda mejor tras el SA. El approach
\texttt{pr\_aislado} ganó 18/23 celdas.

\subsection{Transferencia a las instancias val/egl}
La RTS es la que mejor transfiere de las cinco: su tenure se autoadapta por
diseño, así que los factores calibrados se mantienen sin cambios
(\texttt{factor\_aumento} $= 1.2$, \texttt{factor\_reduccion} $= 0.95$). Solo
se ajusta $p_{inter} = 0.6$ en la clase egl (más vehículos $\Rightarrow$ más
movimientos inter-ruta). Presupuesto uniforme: $10^6$ evaluaciones
($\texttt{iteraciones\_max} = \lfloor 10^6 / \max(20, 2n) \rfloor$) o 300\,s
de pared. Approach: \texttt{pr\_aislado}.
